\section{Governance and Coordination}

The Alliance is a network of independent businesses. The coordination
layer exists not to manage them but to (i) maintain the substrate, (ii)
license the IP, (iii) allocate the capital, (iv) admit new members, and
(v) resolve disputes between members. The discipline is to keep
coordination narrow and leave operations to members.

\subsection{The coordination layer}

The coordination layer sits at \textbf{W3A Foundation}\ (the
Delaware C-corp parent) with \textbf{Hanzo, Inc.}\ as the parallel
licensor on the AI side. Its board includes representatives of the
substrate-developer team, the founding investors, and an independent
director slate.

The coordination layer holds:

\begin{itemize}[leftmargin=2em,itemsep=0.1em]
\item All Tier-2 (LEL) and Tier-3 (LRL--PR) licensing authority.
\item All patent ownership.
\item All trade-secret records.
\item The Alliance brand registry and brand-policy enforcement.
\item The Alliance treasury (Layer 1 capital, \S 6.1).
\item The substrate engineering team.
\item The compliance posture (clean-room protocol, SOC 2 sub-service
organisation responsibility, vendor management).
\end{itemize}

The coordination layer \emph{does not} hold member operations,
member-level customer relationships, member-level product roadmaps, or
member-level hiring decisions. Those remain with each member.

\subsection{Member rights and obligations}

Each member of the Alliance has the following rights:

\begin{itemize}[leftmargin=2em,itemsep=0.1em]
\item Access to the licensed substrate at the executed LEL / LRL--PR /
SCLA terms.
\item Access to the Alliance distribution surface (the cross-member
routing layer, the super-app, the partner ecosystem).
\item Access to the Alliance capital pool at documented terms (Layer 2,
\S 6.2).
\item A seat on the member-coordination forum (quarterly cadence;
information-sharing rather than decision-making).
\item Standard contractual representations and warranties on the
substrate's title, the Alliance's IP, and the partner ecosystem's
performance.
\end{itemize}

And the following obligations:

\begin{itemize}[leftmargin=2em,itemsep=0.1em]
\item Operate under the executed license terms; route through the
Authorized Network; respect the Tier 3 SCLA gate where applicable.
\item Maintain the brand-via-env-override discipline so substrate
upgrades flow cleanly across the network.
\item Maintain SOC 2 sub-service organisation complementary controls
where the member is a regulated counterparty in any jurisdiction.
\item Contribute distribution into the cross-member routing layer at
the standard 50/50-above-documented-costs commercial terms.
\item Respect the clean-room engineering protocol for any work that
extends substrate modules.
\item Report quarterly into the member-coordination forum.
\end{itemize}

\subsection{When the coordination layer intervenes in member operations}

The coordination layer intervenes in member operations only under
narrowly-defined conditions:

\begin{itemize}[leftmargin=2em,itemsep=0.1em]
\item The member is in material breach of its license terms.
\item The member is in material breach of the clean-room protocol.
\item The member's compliance posture has deteriorated to the point of
creating regulatory risk to other members in the same jurisdiction.
\item The member is failing financially and has drawn on member-
stabilisation distribution from the Alliance treasury; the coordination
layer attaches operational conditions to that distribution.
\item A material dispute between two or more members requires
coordination-layer adjudication.
\end{itemize}

In each case the coordination layer's intervention is documented,
time-bounded, and subject to the dispute-resolution mechanism in the
underlying Strategic Partnership Agreement.

\subsection{The PayPal-Mafia parallel --- and where the analogy ends}

The Alliance is deliberately analogous to the PayPal alumni network in
its outcome: a coordinated cohort of people and companies in which one
member's win seeds the next, distribution flows freely across members,
capital recycles into new ventures, and the network as a whole produces
more value than the sum of its parts.

The analogy ends in one important way. The PayPal alumni network was
\emph{discovered ex post}, after the PayPal company itself exited. Its
coordination was informal --- friends helping friends, with no contracted
network economics. The Alliance is the opposite: it is \emph{engineered
ex ante}, with contracted IP licensing, contracted distribution sharing,
contracted capital flows, and contracted governance. The compounding is
guaranteed by the structure rather than left to social capital alone.

\subsection{Exit mechanics}

The Alliance contemplates exits at three levels:

\begin{itemize}[leftmargin=2em,itemsep=0.1em]
\item \textbf{Individual member exit.} A member raises a strategic round
or executes a strategic sale. Alliance retains documented IP-licensing
revenue post-exit; member founders take member-level liquidity. Exit
proceeds flow per the existing cap table.
\item \textbf{Vertical exit.} The Alliance sells or IPOs a complete
vertical (e.g., the entire fintech vertical, or the entire RWA
vertical) to a strategic acquirer or to the public market. The
remaining verticals continue under the Alliance.
\item \textbf{Coordination-layer exit.} The Alliance's parent holding
takes a public listing. The substrate, the IP, the brand registry,
and the network of member relationships all become public-market
assets under a single ticker. Member-level operations continue under
each member's existing structure.
\end{itemize}

The three exit levels are non-mutually-exclusive. A typical multi-year
trajectory contemplates several individual-member exits before any
vertical exit, with the coordination-layer public listing reserved for
year five and beyond.
